\scanpage{33}%
\mspage{18}\space
that it is divisible by $c_4^3-c_6^2$ in $\mathbb{Z}/(2)[c_4,c_6,\Delta]$ (resp. $\mathbb{Z}/(3)[c_4,c_6,\Delta]$). This follows from the structure of the ring of modular forms of characteristic $2$ (resp. $3$).

\uline{Proposition} 7.2. \quad (I) \uline{The ring of modular forms of characteristic} $2$ \uline{is} $\mathbb{Z}/(2)[\tikz[baseline=(coeffa.base)]{\node[inner sep=0pt] (coeffa) {$a_1$};\handwritten{\node[inner sep=0pt,anchor=west,font=\scriptsize] (notea) at ([xshift=1.35em,yshift=-1.05em]coeffa.south west) {the divisor of $a_1$ is smooth over $\mathbb{F}_2$};\draw[thin] (coeffa.south west) to[out=-86,in=192] ([xshift=-.18em]notea.west);}\pgfresetboundingbox\path[use as bounding box] ([yshift=-1.5em]coeffa.south west) rectangle (coeffa.north east);},\Delta]$ , \uline{and} $c_4=a_1^4$ , $c_6=a_1^6$ .

\hspace*{3em}(II) \uline{The ring of modular forms of characteristic} $3$ \uline{is} $\mathbb{Z}/(3)[\tikz[baseline=(coeffb.base)]{\node[inner sep=0pt] (coeffb) {$b_2$};\handwritten{\node[inner sep=0pt,anchor=west,font=\scriptsize] (noteb) at ([xshift=1.35em,yshift=-1.05em]coeffb.south west) {the divisor of $b_2$ is smooth over $\mathbb{F}_2$};\draw[thin] (coeffb.south west) to[out=-86,in=192] ([xshift=-.18em]noteb.west);}\pgfresetboundingbox\path[use as bounding box] ([yshift=-1.5em]coeffb.south west) rectangle (coeffb.north east);},\Delta]$ , and $c_4=b_2^2$ , $c_6=-b_2^3$ .

First of all, it follows from $(1.6)_1$ and $(1.7)_2$ that $a_1$ (resp. $b_2$) is a modular form of characteristic $2$ (resp. $3$) .

In characteristic $2$ , let us first of all look for the modular forms relative only to curves such that $a_1$ is invertible. These curves can be put in the form (cf. (4.\tikz[baseline=(un.base)]{\node[inner sep=0pt] (un) {1};\handwritten{\draw[thin,gray,rotate around={-12:(un.center)}] (un.center) ellipse (0.36em and 1.12em);}}))
\[
 y^2+a_1xy=x^3+a_2x^2+a_6 \qquad \pi=\frac{dx}{a_1x}
\]
and, taking (4.3) into account, one sees that these forms can be written
\[
\begin{aligned}
 &a_1^{-n}\,P(a_1,a_6) \qquad \text{for }P\text{ isobaric, or again}\\[.65em]
 &a_1^{-m}\,Q(a_1,\Delta)\quad\handwritten{,}
\end{aligned}
\]
and this still makes sense for $a_1=0$ , $\Delta\ne0$ if and only if one was in fact dealing with a polyn\tikz[baseline=(oc.base)]{\node[inner sep=0pt] (oc) {o};\handwritten{\draw[thin,gray,rotate around={-14:(oc.center)}] (oc.center) ellipse (0.56em and 1.1em);}}mial in $a_1$ and $\Delta$ .

In characteristic $3$ , reasoning in the same way first on the curves such that $c_4$ is invertible, and noting that $d_6=0$ is then equivalent to $a_4=0$ in (3.1), one sees that every modu form%
% archive sheet number: 32
